\NativeSlot{48}{36}{%
% p062_s01

\ExamQuestion{11（续）}{5}%
最近相距\mbox{（\hspace{1.8em}）}千米。\par
}
\NativeSlot{72}{36}{%
% p062_s02

\ExamQuestion{12.}{5}%
一桶油连桶重\(46.8\)千克，倒出一半油后连桶重\(24.5\)千克，桶重\mbox{（\hspace{1.8em}）}千克，原来有\mbox{（\hspace{1.8em}）}千克油。\par
}
\NativeSlot{123}{43}{%
\ExamGroup{二、选择（将正确答案的序号填在括号里）。{\normalfont\kaishu\fontsize{10.8bp}{14bp}\selectfont （每题2分，共12分）}}

% p062_s03

\ExamQuestion{1.}{5}%
算式\(2\square.4+9.\square6\)的结果可能是\mbox{（\hspace{1.8em}）}。\par
\ExamOptions{A. \(29.50\)}{B. \(29.36\)}{C. \(31.26\)}{}
}
\NativeSlot{204}{56}{%
% p062_s04

\ExamQuestion{2.}{5}%
在右图的加法竖式中，虚线框中的“9”和“5”相加得\mbox{（\hspace{1.8em}）}。\par
\ExamOptions{A. \(14\)个一}{B. \(14\)个十分之一}{C. \(14\)个百分之一}{}
\BookFigure{figures/p062_s04-01.tex}
}
\NativeSlot{303}{57}{%
% p062_s05

\ExamQuestion{3.}{5}%
与\(15.8+9.9\)计算结果不同的是\mbox{（\hspace{1.8em}）}。\par
\ExamOptions{A. \(15.8+9+0.9\)}{B. \(15.8+10+0.1\)}{C. \(15.8+10-0.1\)}{}
}
\NativeSlot{402}{36}{%
% p062_s06

\ExamQuestion{4.}{5}%
两个数相加，如果一个数增加\(3.57\)，另一个数减少\(4.5\)，和\mbox{（\hspace{1.8em}）}。\par
\ExamOptions{A. 减少\(1.07\)}{B. 增加\(0.93\)}{C. 减少\(0.93\)}{}
}
\NativeSlot{452}{36}{%
% p062_s07

\ExamQuestion{5.}{5}%
下面的数最接近9的是\mbox{（\hspace{1.8em}）}。\par
\ExamOptions{A. \(9.02\)}{B. \(8.99\)}{C. \(9.1\)}{}
}
\NativeSlot{501}{87}{%
% p062_s08

\ExamQuestion{6.}{5}%
下列说法中，正确的有\mbox{（\hspace{1.8em}）}个。①两个数的和是小数，这两个数也一定都是小数。②\(6.4+3.6-6.4+3.6=0\)。③计算小数加减法，应把末位对齐。④在减法算式中，被减数和减数同时增加\(2.5\)，差不变。\par
\ExamOptions{A. \(0\)}{B. \(1\)}{C. \(2\)}{}
}
\NativeSlot{650}{64}{%
\ExamGroup{三、计算。{\normalfont\kaishu\fontsize{10.8bp}{14bp}\selectfont （共28分）}}

% p062_s09

\ExamQuestion{1.}{5}%
直接写得数。\par
\noindent \ExamOralCell{3.25+5.5=}\ExamOralCell{1.7+2.5=}\ExamOralCell{6.76-0.6=}\par
\noindent \ExamOralCell{2.6+8=}\ExamOralCell{3.8+0.12=}\ExamOralCell{2.36+1.64=}\par
\noindent \ExamOralCell{4.2-4=}\ExamOralCell{10-4.5=}\par
}
